\documentclass[11pt]{article}

\usepackage[T1]{fontenc}
\usepackage{lmodern}
\usepackage{amsmath,amssymb,amsthm}
\usepackage{booktabs}
\usepackage{microtype}
\usepackage[margin=1in]{geometry}
\usepackage{xcolor}
\usepackage[colorlinks=true,linkcolor=blue,citecolor=blue,
  urlcolor=blue]{hyperref}

\newtheorem{theorem}{Theorem}
\newtheorem{lemma}[theorem]{Lemma}
\newtheorem{proposition}[theorem]{Proposition}
\newcommand{\Rcal}{\mathcal{R}}
\newcommand{\cC}{\mathcal{C}}

\title{The Diameter of the \(K_{3,3}\) Graph Associahedron Is 12:\\
An Exact Computer-Assisted Proof}
\author{Xiaojian Zeng\\
PBC School of Finance, Tsinghua University, 100190\\
\texttt{xiaojianzeng@foxmail.com}}
\date{}

\begin{document}
\maketitle

\begin{abstract}
We determine the diameter of the graph associahedron of the labelled complete
bipartite graph \(K_{3,3}\), under the usual maximal-tubing convention.  An
exact program constructs all 552 graph search trees recursively, constructs
all 1,380 rotations, and runs breadth-first search from every state.  A
structurally separate maximal-clique enumeration reconstructs the same set of
maximal tubings and the same flip graph.  The eccentricities are 10, 11, and
12, with multiplicities 12, 108, and 432, respectively.  Complete exact BFS
records give the upper bound 12 and certify that an explicit pair of tubings
is at distance 12.  Hence the diameter is exactly 12.  The proof uses only
finite integer and set operations; source code and a machine-readable run are
provided for reproduction.
\end{abstract}

\section{Context and result}

Graph associahedra were introduced as polytopes whose face posets are governed
by compatible connected subgraphs~\cite{CarrDevadoss2006}.  Their vertices can
be represented either by maximal tubings or by recursively defined search
trees, and flips of tubings correspond to rotations of search
trees~\cite{CardinalPourninValenciaPabon2022}.  We work throughout with the
labelled graph; in particular, no quotient by automorphisms is taken.

Cardinal, Pournin, and Valencia-Pabon determined the diameters for complete
split graphs and for sufficiently unbalanced complete bipartite graphs.  Their
Question~2 asks for the balanced complete-bipartite regime
\cite{CardinalPourninValenciaPabon2022}.  For the complete split graph
\(SPK_{3,3}\), their Theorem~8 gives diameter 13.  Monotonicity under adding
edges therefore gives
\(\operatorname{diam}\Rcal(K_{3,3})\leq 13\).  On the other hand, the
true-twin edge-deletion inequality of Gargantini, Pastine, and
Torres~\cite[Theorem~4.7]{GargantiniPastineTorres2024} gives
\[
  13-\binom{3}{2}\leq \operatorname{diam}\Rcal(K_{3,3}).
\]
Thus those results give the interval \([10,13]\).  Corollary~4.8 of the latter
paper also specializes to the lower bound 10; the full two-sided interval does
not follow from that corollary alone.

The contribution of this note is the following exact finite determination.

\begin{theorem}\label{thm:main}
Let \(K_{3,3}\) have labelled parts
\(P=\{p_1,p_2,p_3\}\) and \(Q=\{q_1,q_2,q_3\}\).  Let
\(\Rcal(K_{3,3})\) be the graph of all labelled maximal compatible tubings,
with two tubings adjacent when one non-full tube is exchanged.  Then
\[
  \operatorname{diam}\Rcal(K_{3,3})=12.
\]
\end{theorem}

The literature comparison above is deliberately bounded: it identifies the
nearest results used here, but is not a priority claim.

\section{Conventions and finite encodings}

Put \(G=K_{3,3}\).  A \emph{tube} is a nonempty set
\(A\subseteq V(G)\) for which \(G[A]\) is connected.  Distinct tubes \(A,B\)
are compatible when one contains the other, or when they are disjoint and
\(G[A\cup B]\) is disconnected.  A tubing is a pairwise compatible set of
tubes and is maximal with respect to inclusion.  The full tube \(V(G)\) is
compatible with every tube, so it belongs to every maximal tubing.  Retaining
this common tube does not create a flip.

We encode \(p_1,p_2,p_3,q_1,q_2,q_3\) by \(0,1,\ldots,5\), respectively, and
encode a vertex set \(A\) by the bit mask
\[
  m(A)=\sum_{i\in A}2^i\in\{1,\ldots,63\}.
\]
All canonical orderings used by the verifier are lexicographic orderings of
integer tuples of such masks.  A nonempty induced subgraph of \(G\) is
connected precisely when its vertex set is a singleton or meets both parts.
Consequently the number of tubes, including the full tube, is
\begin{equation}\label{eq:tubes}
  6+(2^3-1)^2=55.
\end{equation}

A search tree on a connected induced subgraph \(G[S]\) is defined recursively:
choose a root \(r\in S\), delete \(r\), choose a search tree on every connected
component of \(G[S\setminus\{r\}]\), and attach the roots of those trees to
\(r\).  For a search tree \(T\), let \(D_T(v)\) be the set consisting of \(v\)
and its descendants.  The sets \(D_T(v)\), over all vertices \(v\), form the
maximal tubing corresponding to \(T\)~\cite{CardinalPourninValenciaPabon2022}.

\begin{lemma}\label{lem:recursion}
The recursive construction just described enumerates every search tree on
\(G[S]\) exactly once.
\end{lemma}

\begin{proof}
Induct on \(|S|\).  The root of a search tree is uniquely recoverable.  After
it is deleted, the child subtrees are, by definition, search trees on the
uniquely determined connected components.  The induction hypothesis gives
both completeness and uniqueness within every component, and the Cartesian
product over the components gives all and only the possible choices.  Distinct
roots or distinct component choices cannot yield the same parent vector.
\end{proof}

The implementation also validates every output independently of its route
through the recursion.  For every vertex \(v\), it checks that \(D_T(v)\) is
connected and that
\begin{equation}\label{eq:local-check}
 \{D_T(c):c\text{ is a child of }v\}
 =\cC\bigl(G[D_T(v)\setminus\{v\}]\bigr),
\end{equation}
where \(\cC(H)\) is the set of vertex sets of the connected components of
\(H\).  Thus each stored parent vector is checked against the recursive
definition locally at every node.

\section{Two exhaustive constructions of the flip graph}

The primary construction applies Lemma~\ref{lem:recursion} to the full
six-vertex set.  It obtains 92 search trees for each possible root, hence 552
in total.  It maps each parent vector to the sorted six-tuple
\((m(D_T(v)))_{v\in V(G)}\) and obtains 552 distinct tuples.

The comparison construction begins instead with the 55 masks counted in
\eqref{eq:tubes}.  It constructs the compatibility graph on those masks and
enumerates all its maximal cliques by the Bron--Kerbosch algorithm
\cite{BronKerbosch1973}.  Since a clique in this graph is exactly a tubing and
a maximal clique is exactly a maximal tubing, this is exhaustive directly
from the definitions.  It obtains 552 maximal cliques, each containing six
tubes.  Their complete canonical set is exactly equal to the 552 descendant-set
tuples from the primary construction.

It remains to construct every edge.  In the comparison construction, two
canonical tubings are joined precisely when their symmetric difference has
size two.  Equivalently, for each tubing one deletes each of its five proper
tubes and groups the resulting ridges by equality.  Every ridge has exactly
two owners.  This gives 1,380 undirected edges.

The primary construction obtains edges by rotations without consulting those
ridges.  Suppose that \(t\) is a child of \(s\).  Promote \(t\) above \(s\),
make the old parent of \(s\) the new parent of \(t\), and make \(t\) the new
parent of \(s\).  For each old child \(c\) of \(t\), reattach \(c\) below
\(s\) exactly when
\begin{equation}\label{eq:rotation-rule}
  N_G(s)\cap D_T(c)\neq\varnothing;
\end{equation}
otherwise leave \(c\) below \(t\).

\begin{lemma}\label{lem:rotation}
Rule~\eqref{eq:rotation-rule} gives the search-tree rotation that exchanges the
tube \(D_T(t)\) with one other proper tube.
\end{lemma}

\begin{proof}
Consider the old descendant set of \(s\) after deleting \(t\).  The component
containing \(s\) contains precisely those old child-components of \(t\) that
have an edge to \(s\), together with the unaffected child-components of
\(s\).  An old child-component of \(t\) having no edge to \(s\) is separated
from that component and from the other child-components.  Hence
\eqref{eq:rotation-rule} gives exactly the components required by the recursive
search-tree definition after \(t\) is promoted.  All descendant tubes except
the old \(D_T(t)\) and its replacement are unchanged.
\end{proof}

For every generated state, the verifier applies Lemma~\ref{lem:rotation} to
each of its five parent--child edges.  It validates the resulting parent vector
using \eqref{eq:local-check}, checks that it belongs to the census, and checks
the reverse rotation.  The 2,760 directed rotations give 1,380 undirected
edges, degree five at every state.  The entire edge set is exactly equal to the
edge set from the symmetric-difference construction.

Table~\ref{tab:census} records the resulting finite census.  The equalities in
the last column are explicit full-set comparisons, not comparisons of counts
alone.

\begin{table}[ht]
\centering
\caption{Exact census and cross-checks.}\label{tab:census}
\begin{tabular}{@{}lll@{}}
\toprule
Object & Exact count & Independent check \\
\midrule
Tubes, including the full tube & 55 & formula \eqref{eq:tubes} \\
Search trees / maximal tubings & 552 & identical canonical state sets \\
Directed rotations & 2,760 & inverse checked for each rotation \\
Undirected flip edges & 1,380 & identical canonical edge sets \\
Vertex degrees & \(5\) for all 552 states & both edge constructions \\
\bottomrule
\end{tabular}
\end{table}

\section{All-pairs distances}

On the common graph, the verifier runs ordinary unweighted breadth-first
search from every one of the 552 states.  In each run a vertex receives its
distance when first reached, after all vertices at smaller distance have been
processed.  This standard layer invariant proves that the assigned labels are
shortest-path distances.  Every run reaches all 552 states.  Thus the
computation covers all
\[
  552^2=304{,}704
\]
ordered pairs, rather than a sample.  The eccentricity histogram is
\begin{equation}\label{eq:eccentricities}
  10:12,\qquad 11:108,\qquad 12:432.
\end{equation}
In particular, every distance is at most 12.

For a concrete diametral witness, let \(T_0\) and \(T_{12}\) be the maximal
tubings whose masks are
\[
 T_0=(1,9,11,15,31,63),\qquad
 T_{12}=(32,36,52,60,62,63).
\]
For example, the proper tubes of \(T_0\) are
\[
 \{p_1\}\subset\{p_1,q_1\}\subset
 \{p_1,p_2,q_1\}\subset\{p_1,p_2,p_3,q_1\}
 \subset\{p_1,p_2,p_3,q_1,q_2\},
\]
and the omitted mask 63 is the full tube.  Table~\ref{tab:path} gives a
12-flip path; each row lists the five proper tube masks, with the common full
tube suppressed.  Consecutive rows exchange exactly one tube.

\begin{table}[ht]
\centering
\caption{An explicit 12-flip path from \(T_0\) to \(T_{12}\).}
\label{tab:path}
\begin{tabular}{@{}rl@{}}
\toprule
Step & Proper-tube masks \\
\midrule
0  & \(1,9,11,15,31\) \\
1  & \(1,2,11,15,31\) \\
2  & \(1,2,4,15,31\) \\
3  & \(1,2,4,15,47\) \\
4  & \(1,2,4,39,47\) \\
5  & \(1,2,4,39,55\) \\
6  & \(2,4,38,39,55\) \\
7  & \(2,4,38,54,55\) \\
8  & \(2,4,38,54,62\) \\
9  & \(4,36,38,54,62\) \\
10 & \(4,36,52,54,62\) \\
11 & \(4,36,52,60,62\) \\
12 & \(32,36,52,60,62\) \\
\bottomrule
\end{tabular}
\end{table}

The complete BFS layer record from \(T_0\) places \(T_{12}\) in layer 12 and
in no earlier layer.  Therefore \(d(T_0,T_{12})=12\).  Together with the
all-pairs upper bound following from \eqref{eq:eccentricities}, this proves
Theorem~\ref{thm:main}.

\section{Reproduction, scope, and disclosure}

The ancillary program \texttt{independent\_search\_tree\_verifier.py}
reconstructs the six-vertex graph from its nine individual edges, enumerates
search trees recursively, validates \eqref{eq:local-check}, constructs direct
rotations, and performs every BFS.  Only after its state construction is
complete does it invoke the structurally separate file
\texttt{verify\_k33\_diameter.py} to construct tubes, maximal compatible sets,
and ridge edges.  From the directory containing the ancillary
\texttt{evidence} folder, the complete deterministic run is
\begin{verbatim}
python3 evidence/independent_search_tree_verifier.py \
  --output evidence/independent_search_tree_run.json
\end{verbatim}
The JSON output contains all 552 BFS layer records, the witness path, the
canonical states, and exact comparison results.  Both programs use only the
Python standard library.  All decisive operations are on integers, bit masks,
finite sets, and queues; there is no randomness, floating-point arithmetic,
numerical tolerance, or automorphism reduction.  Assertions for the expected
totals act only after objects have been generated and cannot generate those
objects.

Separating discovery from exact reconstruction and using a second
implementation is also advocated and instantiated in the finite-certificate
workflow of Zeng, Liu, Ratnavelu, and Huat~\cite{ZengEtAl2026Certificates}.
That citation is methodological only: the present theorem is supported by the
definitions, reductions, source code, and complete run described here.

The result is limited to the fixed labelled graph \(K_{3,3}\) and the tubing
and flip conventions in Theorem~\ref{thm:main}.  It makes no assertion about a
formula for the remaining balanced family.  A bounded source search did not
locate another determination of this fixed value, but this is not evidence of
priority.

OpenAI Codex assisted with literature search, proof organization, verifier
development, adversarial audit organization, and initial manuscript
preparation.  A fresh model context separately reconstructed the census and
distances; that was a review pass, not statistical independence or peer
review.  Language-model output is not used as mathematical evidence.  The
proof rests on the human-readable finite reductions and the exact verifiers
accompanying this note.  No proof assistant was used.

\bibliographystyle{plainurl}
\bibliography{references}

\end{document}
